\subsection{Agreement when the specifications are verified}
\label{supp:verified-specification-testing}
The two strategies need not produce the same states. The relevant consistency
property is that a verified specification cannot be falsified by choosing
another legal summary model. Here \emph{verified} means that the complete
modular proof obligations have been discharged in the same interpretation
$\Gamma$: the entry and return contracts, arithmetic and memory safety,
call obligations, and every assertion and invariant selected by the
observation schedule. External functions without bodies contribute explicit
environment assumptions. Proof dependencies and target-layout assumptions
must agree with those used by testing.

Let $A_\ell(s)$ be the assertion established by the verification proof at
program point $\ell$. The proof establishes entry inclusion in $A_{\rm entry}$,
preservation by ordinary program steps, and the selected checks under
$A_\ell$. At a boundary $b$, write $C_b(s,t)$ for the complete relational
postcondition, including old values, the returned value and preserved frames.
Verification of the continuation establishes
\[
 A_\ell(s)\land C_b(s,t)\ \Longrightarrow\ A_{\ell'}(t).
\]
This implication quantifies over \emph{every} legal post-state $t$.
When it returns, the verified body supplies a post-state satisfying $C_b$
in inline mode;
a correct sampler supplies such a state in modular mode.
The same argument covers invariant initialization, preservation and exit,
and continuation after an assertion. The sampler may choose a different
state on every visit; it must preserve the full relational constraints.

\paragraph{From verified assertions to original VCs.}
Trace correctness alone does not establish an entailment over a larger
symbolic source set. For each original record $e$ at point $\ell_e$, the
instance-construction interface therefore requires
\[
  \forall s.\ S_e(s)\Rightarrow A_{\ell_e}(s),
  \qquad
  \forall s.\ A_{\ell_e}(s)\Rightarrow Q_e(s).
\]
The second implication comes from the verified check and faithful resolution;
the first requires that specialization, resource transfer, framing and
normalization retain the verified facts. Together they establish the
\emph{whole} original VC $\phi_e=\forall s.\,S_e(s)\Rightarrow Q_e(s)$,
not merely its truth at the chosen sample.

\begin{theorem}[Agreement for verified specifications]
\label{supp:verified-specification-theorem}
Assume the modular verification obligations above, legal instance inputs,
correct relational sampling, and faithful original-VC construction.
Every finite inline or modular test prefix preserves the verified
assertions and satisfies its selected checks; every generated original VC
is valid. If both tests complete execution and the required coverage, have
nonempty entry domains, and finish soundly deciding every required VC, then
their verdicts agree: both are \pass{}. The claim quantifies over every
legal sequence of modular samples and does not require equal states or
matching execution traces.
\end{theorem}
\begin{proof}
Induct on execution steps. Entry verification establishes the initial
assertion. Ordinary-step verification preserves it. At a contract,
invariant or assertion boundary, legal sampling satisfies the complete
boundary relation; the universally verified continuation implication
establishes the next assertion independently of which legal state was
chosen. The verified assertion entails safety and the selected checks,
including the return condition at a completed exit.
For each original VC, compose source inclusion with its verified target
implication. Sound completed decisions therefore classify all required VCs
as VALID. Nonempty entry and completed coverage then yield \pass{} in both
strategies by the aggregation rule.
\end{proof}

For example, a helper may return 1 while guaranteeing only
$0\le\mathit{ret}\le1$. If the caller's verified continuation accepts this
entire range, inline can produce 1 and modular can sample 0 while both tests
pass. If the caller instead requires exactly 1, the continuation implication
has not been verified against that weak summary. Such a disagreement is
outside the theorem's premise; it does not refute verified-specification
consistency.

The module \texttt{VerifiedSpecificationTesting.v} mechanizes preservation
under legal sampling, the VC implication bridge, and agreement of completed
decisions. Its proof interfaces are explicit; the theorem is not an
end-to-end verification of the C executor or of its solver adapters.
Partial-correctness verification does not guarantee termination or a finite
solver budget. An unfinished execution or proof remains UNKNOWN rather
than contradicting semantic agreement.
The stronger two-way trace correspondence in \texttt{ModeCorrespondence.v}
is retained as supplementary material; it is not required for this theorem.
